\ifdefined\gkver\else\def\gkver{student}\fi
\documentclass[\gkver]{gaokaozhenti}
\begin{document}

\chapter{2006年北京理综}
%% paperType: 理综物理部分

\begin{enumerate}
\item
%% number: 13
%% paperName: 2006年北京理综第 13 题 6 分
%% typeId: 1
%% type: 单选题
%% chapter: 原子和原子核
%% point: 核裂变
%% method: 无
%% score: 6
%% degree: 850
%% duplicateId: 0
%% body:
目前核电站利用的核反应是\xzanswer{A}

\fourchoices[answer=A]
{裂变，核燃料为铀}
{聚变，核燃烧为铀}
{裂变，核燃烧为氘}
{聚变，核燃料为氘}

%% answer: A
%% memo:
\memoanswer{
裂变反应是重核分裂成几个中等质量的原子核的核反应，聚变反应是两个轻核结合成质量较大的原子核的核反应。目前的核电站以 $\ce{^{235}_{92}U}$ 为核燃料，$\ce{^{235}_{92}U}$ 俘获一个中子后发生裂变反应，其核反应方程为 $\ce{^{235}_{92}U + ^{1}_{0}n -> ^{141}_{56}Ba + ^{92}_{36}Kr + 3^{1}_{0}n}$，故 A 正确。
}

\item
%% number: 14
%% paperName: 2006年北京理综第 14 题 6 分
%% typeId: 1
%% type: 单选题
%% chapter: 电场
%% point: 静电
%% method: 无
%% score: 6
%% degree: 650
%% duplicateId: 0
%% body:
使用电的金属球靠近不带电的验电器，验电器的箔片开。下列各图表示验电器上感应电荷的分布情况，正确的是\xzanswer{B}

\fourchoices[answer=B, ispicture=true, h=2.8cm]
{figs/14a.png}
{figs/14b.png}
{figs/14c.png}
{figs/14d.png}

%% answer: B
%% memo:
\memoanswer{
验电器中的自由电子处在带电金属球所形成的电场中，受电场力作用发生定向移动，在验电器的金属球和箔片两端出现等量的异种电荷（感应电荷），当感应电荷产生的电场与带电金属球的电场的矢量和为零时，验电器中的自由电子停止定向移动，达到静电平衡状态，故 A、C 错误；由于电子受电场力方向与电场的方向相反，故若金属球带负电，则验电器中的自由电子将远离金属球，即靠近带电金属球的一端出现的感应电荷为正电荷，而远离带电金属球的一端出现的感应电荷为负电荷，故 B 正确，D 错误。
}

\item
%% number: 15
%% paperName: 2006年北京理综第 15 题 6 分
%% typeId: 1
%% type: 单选题
%% chapter: 气体性质
%% point: 热力学第一定律
%% method: 无
%% score: 6
%% degree: 650
%% duplicateId: 0
%% body:
如图所示，两个相通的容器 P、Q 间装有阀门 K，P 中充满气体，Q 为真空，整个系统与外界没有热交换。打开阀门 K 后，P 中的气体进入 Q 中，最终达到平衡，则\xzanswer{D}

\onepicture[w=5.5cm]{figs/15.png}

\fourchoices[answer=D]
{气体体积膨胀，内能增加}
{气体分子势能减少，内能增加}
{气体分子势能增加，压强可能不变}
{Q 中气体不可能自发地全部退回到 P 中}

%% answer: D
%% memo:
\memoanswer{
“系统与外界没有热交换”表明气体状态的变化过程是绝热过程，即热力学第一定律 $\Delta U = W + Q$ 中 $Q = 0$；由于容器 Q 为真空，故 P 中气体向 Q 中的扩散过程为自由膨胀过程，即 $W = 0$，故气体的内能不变，故 A、B 错误（B 项也可这样判断：由于气体分子间的相互作用力始终表现为引力，故气体膨胀过程分子力做负功，分子势能增大，故 B 错误）；由于气体内能不变而分子势能增大，则气体分子的平均动能减小，故气体温度降低，由 $\frac{pV}{T} = C$ 知气体的压强必定减小，故 C 错误（C 项也可这样判断：由于气体膨胀、分子平均速率减小，故单位时间内气体分子容器壁单位面积的冲量减小，宏观上的气体压强亦减小，故 C 错误）；宏观上的热现象均是不可逆的，故 D 正确。
}

\item
%% number: 16
%% paperName: 2006年北京理综第 16 题 6 分
%% typeId: 1
%% type: 单选题
%% chapter: 光学
%% point: 光的折射
%% method: 无
%% score: 6
%% degree: 650
%% duplicateId: 0
%% body:
水的折射率为 $n$，距水面深 $h$ 处有一个点光源，岸上的人看到水面被该光源照亮的圆形区域的直径为\xzanswer{A}

\fourchoices[answer=A]
{$2h\tan\left(\arcsin\frac{1}{n}\right)$}
{$2h\tan\left(\arcsin n\right)$}
{$2h\tan\left(\arccos\frac{1}{n}\right)$}
{$2h\cot\left(\arccos n\right)$}

%% answer: A
%% memo:
\memoanswer{
当光从水中射向空气时，若入射角大于等于临界角，将在水面处发生全反射现象。如图所示，设光源发出的某条光线恰好在 P 处发生全反射现象，图中 OP 即岸上的人观察到的水面被光源照亮的圆形区域的半径。依临界角的定义有 $\sin C = \frac{1}{n}$，由几何关系有 $R = h\tan C$，解以上两式可得所求圆形区域直径为 $D = 2h\tan\left(\arcsin\frac{1}{n}\right)$，故 A 正确。

\onepicture[w=4.5cm]{figs/16_an.jpg}
}

\item
%% number: 17
%% paperName: 2006年北京理综第 17 题 6 分
%% typeId: 1
%% type: 单选题
%% chapter: 机械振动
%% point: 受迫振动共振
%% method: 无
%% score: 6
%% degree: 650
%% duplicateId: 0
%% body:
某同学看到一只鸟落在树枝上的 P 处，树枝在 $10\Us$ 内上下振动了 6 次，鸟飞走后，他把 $50\Ug$ 的砝码挂在 P 处，发现树枝在 $10\Us$ 内上下振动了 12 次。将 $50\Ug$ 的砝码换成 $500\Ug$ 砝码后，他发现树枝在 $15\Us$ 内上下振动了 6 次，你估计鸟的质量最接近\xzanswer{B}

\onepicture[w=5.5cm]{figs/17.png}

\fourchoices[answer=B]
{$50\Ug$}
{$200\Ug$}
{$500\Ug$}
{$550\Ug$}

%% answer: B
%% memo:
\memoanswer{
由题中所述实验可以看出：在 P 处悬挂物体的质量越大，则树枝振动的频率越低。因此可以断定鸟的质量介于 $50\Ug$ 与 $500\Ug$ 之间，由题给选项可知只有选项 B 可能是正确的。
}

\item
%% number: 18
%% paperName: 2006年北京理综第 18 题 6 分
%% typeId: 1
%% type: 单选题
%% chapter: 曲线运动
%% point: 人造地球卫星
%% method: 无
%% score: 6
%% degree: 450
%% duplicateId: 0
%% body:
一飞船在某行星表面附近沿圆轨道绕该行星飞行。认为行星是密度均匀的球体，要确定该行星的密度，只需要测量\xzanswer{C}

\fourchoices[answer=C]
{飞船的轨道半径}
{飞船的运行速度}
{飞船的运行周期}
{行星的质量}

%% answer: C
%% memo:
\memoanswer{
“飞船在某行星表面附近沿圆轨道绕该行星飞行”，可以认为飞船的轨道半径与行星的半径相等，飞船做圆周运动的向心力由行星对它的万有引力提供，由万有引力定律和牛顿第二定律得 $G\frac{Mm}{R^{2}}=m\left(\frac{2\pi}{T}\right)^{2}R$，由上式可得 $\frac{M}{\frac{4\pi}{3}R^{3}}=\frac{4\pi^{2}}{\frac{4\pi}{3}GT^{2}}$，可得行星的密度 $\rho=\frac{3\pi}{GT^{2}}$，上式表明只要测得飞船运行的周期，即可得到行星的密度，故 C 正确。
}

\item
%% number: 19
%% paperName: 2006年北京理综第 19 题 6 分
%% typeId: 1
%% type: 单选题
%% chapter: 力物体的平衡
%% point: 物体系平衡
%% method: 整体与隔离
%% score: 6
%% degree: 650
%% duplicateId: 0
%% body:
木块 A、B 分别重 $50\UN$ 和 $60\UN$，它们与水平地面之间的动摩擦因数均为 0.25；夹在 A、B 之间的轻弹簧被压缩了 $2\Ucm$，弹簧的劲度系数为 $400\UNm$。系统置于水平地面静止不动。现用 $F = 1\UN$ 的水平拉力作用在木块 B 上。如图所示，力 $F$ 作用后\xzanswer{C}

\onepicture[w=7.0cm]{figs/19.png}

\fourchoices[answer=C]
{木块 A 所受摩擦力大小是 $12.5\UN$}
{木块 A 所受摩擦力大小是 $11.5\UN$}
{木块 B 所受摩擦力大小是 $9\UN$}
{木块 B 所受摩擦力大小是 $7\UN$}

%% answer: C
%% memo:
\memoanswer{
由题给条件知未施加力 $F$ 时，弹簧的弹力大小为 $f = k\cdot\Delta x = 400 \times 0.02\UN = 8\UN$，物块 A 与地面间的滑动摩擦力大小为 $f_{\mu A} = \mu N_{A} = 0.25 \times 50\UN = 12.5\UN$，物块 B 与地面间的滑动摩擦力大小为 $f_{\mu B} = \mu N_{B} = 0.25 \times 60\UN = 15\UN$。施加力 $F$ 后装置仍处于静止状态，B 受地面的摩擦力为 $f_{B}$，A 受地面的摩擦力为 $f_{A}$，由平衡条件得 $F + f = f_{B}$，$f_{A} = f$，代入数据解得 $f_{B} = 9\UN$，$f_{A} = 8\UN$。综合以上分析可知，所给选项中只有 C 项正确。
}

\item
%% number: 20
%% paperName: 2006年北京理综第 20 题 6 分
%% typeId: 1
%% type: 单选题
%% chapter: 磁场
%% point: 带电粒子在磁场中的运动
%% method: 无
%% score: 6
%% degree: 650
%% duplicateId: 0
%% body:
如图所示，匀强磁场的方向垂直纸面向里，一带电微粒从磁场边界 d 点垂直于磁场方向射入，沿曲线 dPa 打到屏 MN 上的 a 点，通过 Pa 段用时为 $t$。若该微粒经过 P 点时，与一个静止的不带电微粒碰撞并结合为一个新微粒，最终打到屏 MN 上。两个微粒所受重力均忽略。新微粒运动的\xzanswer{D}

\onepicture[w=5.0cm]{figs/20.png}

\fourchoices[answer=D]
{轨迹为 Pb，至屏幕的时间将小于 $t$}
{轨迹为 Pc，至屏幕的时间将大于 $t$}
{轨迹为 Pb，至屏幕的时间将等于 $t$}
{轨迹为 Pa，至屏幕的时间将大于 $t$}

%% answer: D
%% memo:
\memoanswer{
令带电微粒的质量为 $m_{1}$，速率为 $v_{1}$，不带电的微粒质量为 $m_{2}$。带电微粒与不带电微粒在 P 处碰撞后的速率为 $v_{2}$，由动量守恒定律得 $m_{1}v_{1} = (m_{1} + m_{2})v_{2}$。由粒子在磁场中做匀速圆周运动的半径 $R = \frac{mv}{qB}$ 可知，碰撞前后微粒做圆周运动的半径未发生变化，即碰后粒子的运动轨迹仍为 Pa；粒子在 Pa 段的飞行时间为 $t = \frac{Pa}{v_{1}}$，$t' = \frac{Pa}{v_{2}}$，因 $v_{2} = \frac{m_{1}v_{1}}{m_{1} + m_{2}} < v_{1}$，所以 $t' > t$，综合可知 D 正确。

\onepicture[w=5.0cm]{figs/20_an.jpg}
}

\item
%% number: 21
%% paperName: 2006年北京理综第 21 题 12 分
%% typeId: 4
%% type: 实验题
%% chapter: 电路
%% point: 电路中的图象
%% method: 图象法
%% score: 12
%% degree: 650
%% duplicateId: 0
%% body:
\begin{enumerate}
	\item
	游标为 20 分度（测量值可准确到 $0.05\Umm$）的卡尺示数如图所示，两测脚间狭缝的宽度为 \tkanswer[1.5cm]{0.15} $\Umm$。用激光照射该狭缝，在屏上出现衍射条纹。如果减小狭缝的宽度，衍射条纹的宽度将变 \tkanswer[1.5cm]{宽}。

	\onepicture[w=7.0cm, align=right]{figs/21.png}

	\item
	某同学用如图所示电路，测绘标有“$3.8\UV$，$0.3\UA$”的小灯泡的灯丝电阻 $R$ 随电压 $U$ 变化的图象：

	\onepicture[w=4.5cm, align=right]{\includesvg{figs/21a}}

	①除了导线和开关外，有以下一些器材可供选择：

	电流表：A$_{1}$（量程 $100\UmA$，内阻约 $2\UO$）；A$_{2}$（量程 $0.6\UA$，内阻约 $0.3\UO$）；

	电压表：V$_{1}$（量程 $5\UV$，内阻约 $5\UkO$）；V$_{2}$（量程 $15\UV$，内阻约 $15\UkO$）；

	滑动变阻器：$R_{1}$（阻值范围 $0\sim10\UO$）；$R_{2}$（阻值范围 $0\sim2\UkO$）；

	电源：$E_{1}$（电动势为 $1.5\UV$，内阻为 $0.2\UO$）；$E_{2}$（电动势为 $4\UV$，内阻约为 $0.04\UO$）。

	为了调节方便，测量准确，实验中应选用电流表 \tkanswer[1.2cm]{A$_{2}$}，电压表 \tkanswer[1.2cm]{V$_{1}$}，滑动变阻器 \tkanswer[1.2cm]{$R_{1}$}，电源 \tkanswer[1.2cm]{$E_{2}$}。（填器材的符号）

	②根据实验数据，计算并描绘出 $R-U$ 的图象如图乙所示。由图象可知，此灯泡在不工作时，灯丝电阻为 \tkanswer[1.2cm]{1.5} $\UO$；当所加电压为 $3.00\UV$ 时，灯丝电阻为 \tkanswer[1.2cm]{11.5} $\UO$，灯泡实际消耗的电功率为 \tkanswer[1.2cm]{0.78} $\UW$。

	\onepicture[w=6.0cm, align=right]{\includesvg{figs/21b}}

	③根据 $R-U$ 图象，可确定小灯泡耗电功率 $P$ 与外加电压 $U$ 的关系。符合该关系的示意图是下列图中的 \tkanswer[1.2cm]{A}。

	\fourchoices[answer=A, ispicture=true, h=2.4cm]
	{figs/21c.png}
	{figs/21d.png}
	{figs/21e.png}
	{figs/21f.png}
\end{enumerate}
%% answer:
\jdanswer{
\begin{enumerate}
	\item
	$0.15$，宽

	\item
	① $E_{2}$，A$_{2}$，V$_{1}$，$R_{1}$；② $1.5$，$11.5$，$0.78$；③ A
\end{enumerate}
}
%% memo:
\memoanswer{
（1）从卡尺的示数可以看出：主尺读数为 $0\Umm$，游标尺读数为 $3 \times 0.05 = 0.15\Umm$，所以所测狭缝宽度为 $0 + 0.15 = 0.15\Umm$；发生明显衍射现象的条件是：障碍物（或孔）的尺寸比波长小或跟波长差不多，本题中若减小狭缝的宽度，衍射现象将更加明显，即衍射条纹将变宽。

（2）①描绘小灯泡的电阻随电压的变化规律，即要求测出小灯泡两端的电压从零到额定电压各个不同状态的电阻阻值，为保证能够测得额定电压下的电阻，则要求电源的电动势略大于额定电压，所以电源应选 $E_{2}$；由小灯泡的标称值知其额定电流为 $0.3\UA$，所以选用的电流表的量程应略大于 $0.3\UA$，故电流表应选 A$_{2}$；为减小读数误差，在电压表量程满足条件的前提下，应尽量选择小量程的电压表，故电压表应选 V$_{1}$；为了调节电压方便，应选择小阻值的滑动变阻器，故滑动变阻器应选 $R_{1}$。

②小灯泡不工作时的电阻，即其两端电压为零时的阻值，由题给 $R-U$ 图象可以看出，当灯泡两端电压为 0 时所对应的电阻为 $1.5\UO$；由图象可以看出，当灯泡两端的电压为 $3.00\UV$ 时，所对应的电阻阻值为 $11.5\UO$；灯泡的功率为 $P=\frac{U^{2}}{R}=\frac{3.00^{2}}{11.5}\UW=0.78\UW$。

③当灯泡两端的电压为 0 时，灯泡的功率也为零，故 B、D 错误；由 $P = \frac{U^{2}}{R}$ 可在图象中选取几个特殊点，分别计算在不同电压下灯泡的功率，定性地作出 $P-U$ 图象，可得图象 A 正确。
}

\item
%% number: 22
%% paperName: 2006年北京理综第 22 题 16 分
%% typeId: 6
%% type: 计算题
%% chapter: 功和能
%% point: 动能定理
%% method: 无
%% score: 16
%% degree: 650
%% duplicateId: 0
%% body:
下图是简化后的跳台滑雪的雪道示意图。整个雪道由倾斜的助滑雪道 AB 和着陆雪道 DE，以及水平的起跳平台 CD 组成，AB 与 CD 圆滑连接。运动员从助滑雪道 AB 上由静止开始，在重力作用下，滑到 D 点水平飞出，不计飞行中的空气阻力，经 $2\Us$ 在水平方向飞行了 $60\Um$，落在着陆雪道 DE 上，已知从 B 点到 D 点运动员的速度大小不变。（$g$ 取 $10\Umsq$）求：

\begin{enumerate}
	\item
	运动员在 AB 段下滑到 B 点的速度大小。
	\item
	若不计阻力，运动员在 AB 段下滑过程中下降的高度。
	\item
	若运动员的质量为 $60\Ukg$，在 AB 段下降的实际高度是 $50\Um$，此过程中他克服阻力所做的功。
\end{enumerate}

\onepicture[w=6.0cm, align=right]{figs/22.png}
%% answer:
\jdanswer{
\begin{enumerate}
	\item
	$v = 30\Ums$
	\item
	$h = 45\Um$
	\item
	$W_{f} = 3000\UJ$
\end{enumerate}
}
%% memo:
\memoanswer{
（1）运动员从 D 点飞出时的速度 $v=\frac{s_{x}}{t}=30\Ums$。依题意，下滑到助滑雪道末端 B 点的速度大小是 $30\Ums$。

（2）在下滑过程中机械能守恒，有 $mgh=\frac{1}{2}mv^{2}$，下降的高度 $h=\frac{v^{2}}{2g}=45\Um$。

（3）由动能定理得 $mgh - W_{f} = \frac{1}{2}mv^{2}$，解得运动员克服阻力做功 $W_{f}=mgH-\frac{1}{2}mv^{2}=3000\UJ$。
}

\item
%% number: 23
%% paperName: 2006年北京理综第 23 题 18 分
%% typeId: 6
%% type: 计算题
%% chapter: 电场
%% point: 带电粒子的直线运动
%% method: 无
%% score: 18
%% degree: 450
%% duplicateId: 0
%% body:
如图 1 所示，真空中相距 $d = 5\Ucm$ 的两块平行金属板 A、B 与电源连接（图中未画出），其中 B 板接地（电势为零），A 板电势变化的规律如图 2 所示。将一个质量 $m = 2.0\times10^{-23}\Ukg$，电量 $q = +1.6\times10^{-19}\UC$ 的带电粒子从紧临 B 板处释放，不计重力。求：

\begin{enumerate}
	\item
	在 $t = 0$ 时刻释放该带电粒子，释放瞬间粒子加速度的大小；
	\item
	若 A 板电势变化周期 $T = 1.0\times10^{-5}\Us$，在 $t = 0$ 时将带电粒子从紧临 B 板处无初速释放，粒子到达 A 板时动量的大小；
	\item
	A 板电势变化频率多大时，在 $t = \frac{T}{4}$ 到 $t = \frac{T}{2}$ 时间内从紧临 B 板处无初速释放该带电粒子，粒子不能到达 A 板。
\end{enumerate}

\onepicture[w=7.5cm, align=right]{figs/23.png}
%% answer:
\jdanswer{
\begin{enumerate}
	\item
	$a = 4\times10^{9}\Umsq$
	\item
	$p = 4.0\times10^{-23}\Ukgms$
	\item
	$f > 5\sqrt{2}\times10^{4}\UHz$
\end{enumerate}
}
%% memo:
\memoanswer{
（1）金属板间的电场强度为 $E=\frac{U}{d}$，带电粒子所受电场力 $F=qE=\frac{qU}{d}$，由牛顿第二定律得 $F=ma$，解得 $a=\frac{qU}{md}=4\times10^{9}\Umsq$。

（2）粒子在 $0\sim\frac{T}{2}$ 时间内走过的距离为 $\frac{1}{2}a\left(\frac{T}{2}\right)^{2}=5.0\times10^{-2}\Um$，故带电粒子在 $t=\frac{T}{2}$ 时恰好到达 A 板。根据动量定理，此时粒子动量 $p = Ft = 4.0\times10^{-23}\Ukgms$。

（3）带电粒子在 $t=\frac{T}{4}\sim t=\frac{T}{2}$ 向 A 板做匀加速运动，在 $t=\frac{T}{2}\sim t=\frac{3T}{4}$ 向 A 板做匀减速运动，速度为零后将返回，粒子向 A 板运动可能的最大位移 $s=2\times\frac{1}{2}a\left(\frac{T}{4}\right)^{2}=\frac{1}{16}aT^{2}$。要求粒子不能到达 A 板，有 $s < d$。由 $f=\frac{1}{T}$，电势变化频率应满足 $f>\sqrt{\frac{a}{16d}}=5\sqrt{2}\times10^{4}\UHz$。
}

\item
%% number: 24
%% paperName: 2006年北京理综第 24 题 20 分
%% typeId: 6
%% type: 计算题
%% chapter: 磁场
%% point: 安培力
%% method: 无
%% score: 20
%% degree: 450
%% duplicateId: 0
%% body:
磁流体推进船的动力来源于电流与磁场间的相互作用。图（a）是在平静海面上某实验船的示意图，磁流体推进器由磁体、电极和矩形通道（简称通道）组成。如图（b）所示，通道尺寸 $a = 2.0\Um$、$b = 0.15\Um$、$c = 0.10\Um$。工作时，在通道内沿 $z$ 轴正方向加 $B = 8.0\UT$ 的匀强磁场；沿 $x$ 轴负方向加匀强电场，使两金属板间的电压 $U = 99.6\UV$；海水沿 $y$ 轴方向流过通道。已知海水的电阻率 $\rho = 0.20\UOm$。

\begin{enumerate}
	\item
	船静止时，求电源接通瞬间推进器对海水推力的大小和方向；
	\item
	船以 $v_{s} = 5.0\Ums$ 的速度匀速前进。若以船为参照物，海水以 $5.0\Ums$ 的速率涌入进水口，由于通道的截面积小于进水口的截面积，在通道内海水速率增加到 $v_{d} = 8.0\Ums$。求此时两金属板间的感应电动势 $U_{\text{感}}$；
	\item
	船行驶时，通道中海水两侧的电压按 $U' = U - U_{\text{感}}$ 计算，海水受到电磁力的 $80\%$ 可以转化为对船的推力。当船以 $v_{s} = 5.0\Ums$ 的速度匀速前进时，求海水推力的功率。
\end{enumerate}

\onepicture[w=8.0cm, align=right]{\includesvg{figs/24}}
%% answer:
\jdanswer{
\begin{enumerate}
	\item
	$F_{1} = 796.8\UN$，方向沿 $y$ 轴正方向（向右）
	\item
	$U_{\text{感}} = 9.6\UV$
	\item
	$P = 2880\UW$
\end{enumerate}
}
%% memo:
\memoanswer{
（1）根据安培力公式知推进器对海水的推力为 $F_{1}=I_{1}Bb$，其中 $I_{1}=\frac{U}{R}$，又 $R=\rho\frac{b}{ac}$，联立解得 $F_{1}=\frac{U}{R}Bb=\frac{Uac}{\rho}B=796.8\UN$。对海水推力的方向沿 $y$ 轴正方向（向右）。

（2）$U_{\text{感}} = Bv_{d}b = 9.6\UV$。

（3）根据欧姆定律 $I_{2}=\frac{U'}{R}=\frac{(U-Bv_{d}b)ac}{\rho b}=600\UA$，安培推力 $F_{2}=I_{2}Bb=720\UN$，对船的推力 $F=80\%F_{2}=576\UN$，推力的功率 $P=Fv_{s}=80\%F_{2}v_{s}=2880\UW$。
}

\end{enumerate}

\end{document}
